\chapter{Einleitung}\label{ch:einleitung}


\section{Problemstellung und Motivation}\label{sec:problemstellung}

\gls{erp}-Systeme sind das zentrale Anwendungssystem für den betrieblichen Alltag. Sie umfassen nach Gronau die Verwaltung aller zur Durchführung von Geschäftsprozessen notwendigen Informationen über die Ressourcen Material, Personal, Kapazitäten, Finanzen und Informationen und setzen dafür eine gemeinsame Datenhaltung voraus\footcite[Vgl.][S.~4]{gronau2014erp}. Auf der operativen Ebene führen sie die täglichen Routinetransaktionen des Geschäftsbetriebs aus und zeichnen diese auf\footcite[Vgl.][S.~125 ff.]{vanderaalst2016processmining}. Jeder Beleg, jede Buchung und jede Änderung hinterlässt dabei in der zentralen Datenbank eine digitale Spur mit Zeitstempel. Process Mining macht diese Spuren nutzbar und analysiert Geschäftsprozesse nicht so, wie sie dokumentiert sind, sondern so, wie sie tatsächlich ablaufen\footcites(Vgl.)()[S.~169 f.]{vanderaalst2012manifesto}[S.~13]{adolfsson2026o2c}.

Untersuchungsgegenstand ist der \gls{o2c} im Vertriebsmodul \gls{sd} von \gls{sap} \gls{s/4hana}. Er umfasst die Prozesskette von der Auftragsaufnahme über Verfügbarkeitsprüfung, Kommissionierung und Versand bis zur Fakturierung und zum Zahlungseingang\footcite[Vgl.][]{SAP_Verkaufsanalyse} und verläuft als erlöswirksamer End-to-End-Prozess abteilungsübergreifend durch Vertrieb, Lager und Finanzbuchhaltung\footcite[Vgl.][S.~63 ff.]{scheibler2021vertrieb}. Sein hohes Belegvolumen, die unmittelbare Wirkung auf Kundenzufriedenheit und Liquidität sowie die vollständige Abbildung im \gls{erp}-System eignen ihn gut für eine Process-Mining-Analyse. Zugleich gilt der \gls{o2c}-Prozess als einer der am besten standardisierten und zugleich am häufigsten durch gewachsene Sonderregelungen überformten Prozesse, sodass sich an ihm besonders deutlich zeigt, wie weit dokumentierter Soll-Ablauf und tatsächliche Ausführung auseinanderfallen können.

In der beruflichen Praxis der Auftragsabwicklung zeigen sich regelmäßig die typischen Symptome ineffizienter \gls{o2c}-Prozesse: Aufträge, die in Kredit- oder Liefersperren hängen bleiben, Klärungsschleifen zwischen Vertrieb und Versand nach Auftragsänderungen sowie Rechnungen, die das Haus erst Tage nach dem Warenausgang verlassen. Auch die Fachliteratur zu \gls{sap} \gls{sd} benennt manuelle Dateneingabe, fragmentierte Informationssysteme und verzögerte Auftragsabwicklung als wiederkehrende Schwachstellen\footcite[Vgl.][S.~109 f.]{nadukuru2024o2c}. Dieselbe Literatur identifiziert Automatisierung, Datenintegration und fortgeschrittene Analytik als die zentralen Hebel, um diese Schwächen zu beheben, und ordnet damit die datengetriebene Prozessanalyse in eine breitere Optimierungsagenda ein\footcite[Vgl.][S.~110 ff.]{nadukuru2024o2c}. Versuche, solche Ursachen in Prozessworkshops zu erheben, bleiben erfahrungsgemäß unbefriedigend, weil die Beteiligten denselben Prozess unterschiedlich beschreiben und die Prozessdokumentation häufig unvollständig ist\footcite[Vgl.][S.~50 und 62]{adolfsson2026o2c}. Process Mining stützt die Analyse stattdessen auf Daten: Die Ergebnisse sind faktenbasiert und intersubjektiv nachvollziehbar, weil sie unmittelbar aus den Aufzeichnungen des Systems gewonnen werden\footcite[Vgl.][S.~169 ff.]{vanderaalst2012manifesto}.

Bewusst wird kein reales Unternehmen untersucht, sondern ein offen dokumentierter Datensatz. Die Datenstrukturen und Algorithmen des Process Mining lassen sich an einem reproduzierbaren Event-Log präziser untersuchen als an vertraulichen Unternehmensdaten, deren Freigabe zudem datenschutzrechtliche und arbeitsrechtliche Fragen aufwerfen würde. Die Datengrundlage wird in Kapitel~\ref{sec:datengrundlage} eingeordnet und begründet.

\section{Zielsetzung und Forschungsfragen}\label{sec:zielsetzung}

Ziel der Arbeit ist es, zu untersuchen, wie Prozessineffizienzen in \gls{erp}-Systemen mittels Process Mining identifiziert werden können. Dazu werden drei Leitfragen verfolgt:
\begin{enumerate}
    \item Welche Datenstrukturen eines \gls{sap}-Systems liefern die Rohdaten für ein Process-Mining-taugliches Ereignisprotokoll (Event-Log), und wie wird dieses daraus gewonnen?
    \item Welche Algorithmen stehen für die Prozessentdeckung und den Soll-Ist-Abgleich zur Verfügung, und wie schlagen sie sich im praktischen Vergleich auf einem realistischen \gls{o2c}-Log?
    \item Welche Prozessineffizienzen lassen sich in dem analysierten Log quantifizieren, und mit welchen Anpassungen im \gls{sap}-System ließen sie sich beheben?
\end{enumerate}
Da der verwendete Log synthetisch erzeugt wird, die injizierten Störungsmuster also vorab bekannt sind, dient die Analyse zugleich als Methodenvalidierung. Ein taugliches Analyseverfahren muss die bekannte Grundwahrheit aus den Daten zurückgewinnen. Methodisch folgt die Arbeit damit dem Grundgedanken der gestaltungsorientierten Wirtschaftsinformatik, in der ein Artefakt, hier die Extraktions- und Analysemethode, konstruiert und anschließend an einem nachprüfbaren Maßstab evaluiert wird. Sie liefert damit nicht nur eine Fallstudie, sondern einen überprüfbaren Nachweis, dass sich aus den Standard-Datenstrukturen eines \gls{sap}-Systems belastbare Prozessdiagnosen gewinnen lassen. Der Anspruch der Arbeit ist damit ausdrücklich ein methodischer: Nicht die Behebung eines konkreten Einzelfalls steht im Vordergrund, sondern der nachvollziehbare Nachweis eines übertragbaren, beratungsorientierten Vorgehens.

\section{Vorgehensweise und Aufbau der Arbeit}\label{sec:vorgehensweise}

Die Arbeit folgt der funktionalen Gliederung in Einleitung, Grundlagen, Umsetzung und Schluss. Kapitel~\ref{ch:grundlagen} legt die Grundlagen: den \gls{o2c}-Prozess mit seinen typischen Schwachstellen, seinen Aufbau im Modul \gls{sd} mit Strukturen, Stamm- und Bewegungsdaten sowie die Konzepte und Algorithmen des Process Mining. Kapitel~\ref{ch:umsetzung} bildet die Umsetzung und damit den Kern der Arbeit: Aus den \gls{sap}-Datenstrukturen wird ein Event-Log konstruiert, auf dem Discovery- und Conformance-Verfahren ausgeführt und die Ineffizienzen quantifiziert werden. Kapitel~\ref{ch:empfehlungen} leitet daraus beratungsorientierte Handlungsempfehlungen einschließlich einer Implementierungsstrategie ab. Kapitel~\ref{ch:fazit} fasst die Ergebnisse entlang der drei Forschungsfragen zusammen und gibt einen Ausblick.
